% This is the "preamble" of the document. This is where the format options get set.
% Pro-tip: things following the % mark will not be compiled by LaTeX. I'll be using them extensively to explain things as we go.
% Note: not to scare you off of LaTeX, but it's normal to have problems. And ya girl has been having some. I've included the copyright info at the bottom of the document from the guy who wrote this package, because his documentation doesn't entirely match how it's actually used. So this is a combination of his working preamble along with my added commentary or explanation.
%%


\documentclass[stu,12pt,floatsintext]{apa7}

\usepackage[american]{babel}
\usepackage{authblk}
\usepackage{csquotes} % One of the things you learn about LaTeX is at some level, it's like magic. The references weren't printing as they should without this line, and the guy who wrote the package included it, so here it is. Because LaTeX reasons.
\usepackage[style=apa,sortcites=true,sorting=nyt,backend=biber]{biblatex}
% biblatex: loads the package that will handle the bibliographic info. Other option is natbib, which allows for more customization
% - style=apa: sets the reference format to use apa (albeit the 6th edition)
\DeclareLanguageMapping{american}{american-apa} % Gotta make sure we're patriotic up in here. Seriously, though, there can be local variants to how citations are handled, this sets it to the American idiosyncrasies
\addbibresource{references.bib} % This is the companion file to the main one you're writing. It contains all of the bibliographic info for your references. It's a little bit of a pain to get used to, but once you do, it's the best. Especially if you recycle references between papers. You only have to get the pieces in the holes once.`

\usepackage{comment}
\usepackage{siunitx}
\sisetup{round-mode=places,round-precision=2}\usepackage[T1]{fontenc}
\usepackage{mathptmx} % This is the Times New Roman font, which was the norm back in my day. If you'd like to use a different font, the options are laid out here: https://www.overleaf.com/learn/latex/Font_typefaces
% Alternately, you can comment out or delete these two commands and just use the Overleaf default font. So many choices!


% Title page stuff
\title{Predicting a Successful Climb of Mount Everest} % The big, long version of the title for the title page
\shorttitle{What Defines a Successful Climb of Mount Everest} % The short title for the header
\authorsnames{Jacob Tassos,Mostafa Zamani Turk,Marston S. Ward}
\duedate{June 23, 2025}
% \date{January 17, 2024} The student version doesn't use the \date command, for whatever reason
\affiliation{University of San Diego}
\course{MS-AAI 500-01} % LaTeX gets annoyed (i.e., throws a grumble-error) if this is blank, so I put something here. However, if your instructor will mark you off for this being on the title page, you can leave this entry blank (delete the PSY 4321, but leave the command), and just make peace with the error that will happen. It won't break the document.
\professor{Professor Leonid Shpaner M.S.}  % Same situation as for the course info. Some instructors want this, some absolutely don't and will take off points. So do what you gotta.

\abstract{Abstract to be written later}

%\keywords{APA style, demonstration} % If you need to have keywords for your paper, delete the % at the start of this line

\begin{document}
\maketitle % This tells LaTeX to make the title page

% \section{Introduction} This command is commented out, because I was taught it was redundant to have the paper's title and introduction together. If your instructor wants it to say "Introduction", delete the % at the start

\section{Introduction}
Mount Everest is the highest land surface point on the planet, when this height is defined as height above mean sea level. At \si{8,848~\km} (\si{29,031.7~feet}), it towers above it's closest rival by more than \si{200~\km} (See list of highest mountains on Earth: \url{https://en.wikipedia.org/wiki/List_of_highest_mountains_on_Earth}). While there are more than 100 peaks on the planet with peaks above \si{7~\km}, none captures the human imagination as Everest. Humans cannot live at such high elevations far above the tree line, where the air is rarefied and every function of the human body is challenged \cite{goos_how_2018}.

Despite the extreme environment of peaks about 7,000 km, the human desire to overcome a challenge even at extremely personal risk remains alive and well. Mountains are not the only extreme environment on the planet that humans have tried to reach. The deepest parts of the oceans are also areas we have successfully traversed, however, the costs of specialized equipment and availability of aid in achieving the goal of reaching these extreme environment has made Everest an attractive place because it became attainable and relatively inexpensive.

For centuries now, climbing the worlds highest mountains has been enabled and facilitated by modern technology and, in the case of Everest, Nepal tourist industry built around climbing the mountain. These two factors have made one of the world's deadliest mountains accessible to even more adventurers, and many would say, amateurs. In the early days of climbing teams climb to be the first to reach the highest peak, now anyone with enough resources can travel to, especially Everest and attempt to summit. But far from all who do, make it, or even make it back in one piece - if they make it back at all.

Everest and K2 (second peak in Karakoram) have over the years competed for the title: deadliest mountain. Many lives have been lost due a range of factors ranging from physiology, mindset, preparedness, ego, stupidity and natural disasters such as avalanches and inclement weather. Many studies and documentaries have been done on mountain climbers and their plight to summit and return - alive. But sadly, the history of climbing is replete with tragedies. In the light of this some articles have tried to analyze and suggest ways to improve your chances of a successful summit. In this paper we define that success as one where the climber summits and returns under their own strength, and or course, alive! For example, \textcite{huey_mountaineering_2001} analyzed a research program that attempted to predict (similar to this paper) what factors influence the success ascent or death rates. In this project we do not try to predict death rates as not all failures to summit is due to the mountaineer dying. However, the paper drew no conclusions on the matter. \textcite{szymczak_comparison_2021} examined the weather climber experienced on Everest and K2 in both the climbing season (July) and mid winter season. They study found that more extreme weather in both seasons than on K2, which is \si{8^{\circ}} latitude farther north.


\section{EDA}

\begin{table}
    \centering
    \caption{Results of Exploratory Data Analysis}
    \begin{tabular}{lcccc}
        \hline
         Variable Name & Total Unique Values & Mean ($\mu$) & STD ($\sigma$) & Continuous/Discrete \\
         \hline
         peak$\_$name & 385 & - & - & D\\
         nationality & 97 & - & - & D\\
         year & 93 & 2001.53 & 14.76 & D\\
         season & 5 & - & - & D \\
         total$\_$days & 93 & 20.28 & 17.66 & C\\
         exp$\_$result & 15 & - & - & D\\
         max$\_$elev$\_$reached & 696 & 7,106.91 & 1,763.61 & C\\
         mbrs$\_$deaths & 8 & \num{0.075} & \num{0.374835} & C\\
         is$\_$o2$\_$not$\_$used & 2 & \num{0.70640} & \num{0.455431} & D\\
         is$\_$o2$\_$climbing & 2 & \num{0.249190} & \num{0.432565} & D\\
         is$\_$o2$\_$descent & 2 & \num{0.013913} & \num{0.117134} & D\\
         is$\_$o2$\_$medical & 2 & \num{0.030684} & \num{0.172469} & D\\
         is$\_$o2$\_$unkwn & 2 & \num{0.002287} & \num{0.047770} & D\\
         had$\_$o2 & \num{0.086621} & 2 & \num{0.281292} & D\\
         accidents & 2691 & - & - & D\\
         \hline
    \end{tabular}
    \label{tab:eda}
\end{table}


\begin{figure}
    \centering
    \includegraphics[width=0.75\linewidth]{sampleFig.png}
    \caption{The mean vibes for each word type and counterbalanced condition order. Error bars represent one standard error of the mean.}
    \label{fig:OverallEffect}
\end{figure}


\section{Null Hypothesis}
Participant information goes here.

\section{Materials}

Words, along with a sample table (Table~\ref{tab:table_words}).

% There's lots of little components to the table. For the most part you can just copy it, or build your own with Overleaf's table wizard up top.
% I'll mention that the & character separates items into different columns on a row, and the \\ ends that line. \hline generates a line that matches the width of the table.

\section{Model selection}

Some words about the measures used.

\section{Features and Ranking}
Some more words


\section{Results}

Description of the procedure.

\section{Conclusion}
Placeholder

\printbibliography

\end{document}




\begin{comment}
Based on the research by (Fontanarosa et al., 2000), the use of supplemental oxygen has been shown to significantly reduce the risk of death among mountaineers descending from the summits of Everest and K2. Their findings highlight that climbers who forgo supplemental oxygen face substantially higher mortality rates, particularly on K2. This underscores the critical role of oxygen support in mitigating physiological stress and enhancing survival in extreme high-altitude environments. Since Mount Everest became accessible from both Nepal and Tibet, its climbing history has been shaped by changing geopolitical boundaries and international expedition dynamics. This paper (Huey & Salisbury, 2003) examines patterns of success and mortality across thousands of Everest climbers, revealing how route choice, nationality, and historical context influence mountaineering outcomes. A detailed analysis of 86 years of expedition data reveals that late summit times and symptoms such as fatigue and cognitive impairment are strong indicators of potential mortality at extreme altitudes. These findings emphasize the critical importance of timely decision-making and physical readiness in high-altitude mountaineering (Firth et al., 2008). (Westhoff et al., 2012) found that individual experience or participation in traditional (non-commercial) expeditions does not significantly improve survival odds. Instead, their study highlights that broader, collective progress—such as innovations in equipment, logistics, and shared knowledge—plays a more critical role in reducing fatalities over time. A study by (Anicich et al., 2015) revealed that cultural values related to hierarchy can both enhance and endanger expedition outcomes in the Himalayas. While hierarchical teams were more likely to reach the summit, they also experienced higher death rates—highlighting the complex trade-off between team coordination and psychological safety in high-risk group settings. (Gugglberger, 2018) explored the evolving role of women in Himalayan mountaineering, emphasizing not only their increasing participation but also the unique risks and challenges they face. The analysis sheds light on how gender dynamics intersect with survival outcomes in high-altitude expeditions. A study by (Huey et al., 2020) analyzed nearly 6,000 first-time climbers over two time periods and found that while probabilities of summiting have increased—particularly among women and older climbers—death rates have remained largely unchanged, with age emerging as a significant risk factor. The 2019 Everest Expedition, led by National Geographic and Rolex, marked the most comprehensive scientific study ever conducted on the mountain. The research revealed emerging risks to both climbers and local communities driven by environmental and human-induced changes, highlighting the urgent need to understand and mitigate evolving threats in the Everest and Khumbu region (Miner et al., 2020). (Krishnagopal, 2021) employs a novel multiscale network approach to analyze how both individual traits—such as age, gender, and experience—and expedition-wide factors influence mountaineering success. The findings highlight that climbing with familiar teammates and factors like youth and oxygen use significantly improve success rates, while expedition size and duration also play crucial roles. Mountaineering uniquely combines individual skill and endurance with the critical influence of social dynamics among climbers. Recent research (Krishnagopal, 2022) demonstrates that the structure and strength of relationships within climbing teams significantly affect cooperation levels, summit success, and even death rates, with both individual traits and expedition-wide factors playing important roles in these outcomes. The high-altitude environment presents significant health risks to mountaineers, including cold injuries such as frostbite and the potential for mortality. Recent reviews (Kriemler et al., 2023) suggest that female mountaineers may experience a lower risk of death compared to their male counterparts, though data on sex differences in frostbite remain inconclusive and warrant further investigation.
\end{comment}

